\documentclass[twoside,11pt]{article}
\usepackage{jair, theapa, rawfonts}
\usepackage[dvipsnames]{xcolor}
\usepackage{amsmath,amssymb,amsthm}
\usepackage[ruled,linesnumbered]{algorithm2e}
\usepackage{thmtools} 
\usepackage{booktabs}
\usepackage{paralist}
\usepackage{subcaption}
\usepackage{graphicx}
\usepackage{adjustbox}
\usepackage{multirow,array}
\usepackage[normalem]{ulem}
\usepackage[capitalise]{cleveref}
\usepackage{cancel}
%\definecolor{vuborange}{rgb}{1.0,0.40,0.0}



%%%%%%
\newcommand\citet[1]{\citeauthor{#1}~[\citeyear{#1}]}

\newcommand\m[1]{\mathcal{#1}}



%% Theorems/Theory
\newtheorem{thm}{Theorem}
\newtheorem{definition}[thm]{Definition}
\newtheorem{prop}{Property}
\newtheorem{property}[prop]{Property}
\newtheorem{lem}{Lemma}
\newtheorem{lemma}[lem]{Lemma}
\newtheorem{propo}{Proposition}
\newtheorem{proposition}[propo]{Proposition}
\newtheorem{ex}{Example}
\newtheorem{example}[ex]{Example}
\newtheorem{obs}{Observation}
\newtheorem{observation}[obs]{Observation}
\newtheorem{theorem}[thm]{Theorem}


%% Comments
\newcommand\comment[1]{\marginpar{\tiny #1}}
\renewcommand\comment[1]{#1}
\newcommand{\tias}[1]{{\comment{\color{blue}Tias:#1}}}
\newcommand{\emilio}[1]{{\comment{\color{red}Emilio: #1}}}
\newcommand{\bart}[1]{{\comment{\color{green}Bart: #1}}}
\newcommand{\new}[1]{{\comment{{\color{blue} #1}}}}

%% Factoring
\newcommand\setstohit{\ensuremath{\m{H} }\xspace}
\newcommand\F{\ensuremath{\m{F} }\xspace}
\newcommand\subsetT{\ensuremath{\m{S} }\xspace}
\newcommand\bestsubsetT{\ensuremath{\m{S}_{best}}\xspace}
\newcommand\corrsetT{\ensuremath{\m{K} }\xspace}
\newcommand\setstohitall{\ensuremath{\m{H}_\mathit{all} }\xspace}
\newcommand\Iend{\ensuremath{I_\mathit{end} }\xspace}
\newcommand\formula{\ensuremath{\m{F} }\xspace}
\newcommand\formulac{\ensuremath{\m{C} }\xspace}
\newcommand\formulag{\ensuremath{\m{G} }\xspace}
\newcommand\mm[1]{\ensuremath{#1}\xspace}
\newcommand\nat{\mm{\mathbb{N}}}
\newcommand\ltrue{\mm{\textbf{t}}}
\newcommand\lfalse{\mm{\textbf{f}}}
\newcommand\uservars{\ensuremath{\m{U} }\xspace}
\newcommand\call[1]{\mm{\textsc{#1}}}
\newcommand\cohs{\mm{\call{CondOptHittingSet}}}
\newcommand\corrsubsets{\mm{\call{CorrSubsets}}}
\newcommand\corrsubsetsgreedy{\mm{\call{CorrSubsetsGreedy}}}
\newcommand\disjcorr{\mm{\call{DisjCorrSets}}}
\newcommand\disjcorrincr{\mm{\call{DisjCorrSetsIncremental}}}
\newcommand\disjcorrconstrained{\mm{\call{DisjCorrSetsConstrained}}}
\newcommand\sat{\texttt{{SAT}}\xspace}
\newcommand\grow{\mm{\call{Grow}}}
\newcommand\omus{\mm{\call{OUS}}}
\newcommand\comus{\mm{\call{OCUS}}}
\newcommand\comusdisj{\mm{\call{OCUSConstrainedDisjMCSes}}}
\newcommand\disjmcs{\mm{\call{DisjointMCSes}\xspace}}
\newcommand\disjmcsconstrained{\mm{\call{DisjointMCSesConstrained}\xspace}}
\newcommand\omusinc{\mm{\call{OUS-Inc}}}
\newcommand\hitsetbased{hitting set--based\xspace} %en-dash!
\newcommand\satsets{\mm{\mathbf{SSs}}}
\newcommand\fall{\mm{\formula_{\mathit{all}}}}
\newcommand\algemilio[1]{\emilio{#1}\;}
\newcommand\UB{ub}
\newcommand\LB{lb}
\newcommand\maxsat{MaxSAT\xspace}
\newcommand\voc{\ensuremath{\Sigma}\xspace}
\newcommand\negset[1]{\mm{\overline{#1}}}

\newcommand{\algorithmicbreak}{\textbf{break}}
%%%%%
\newcommand\muses[1]{\ensuremath{\mathit{MUSs}(#1)}\xspace}
\newcommand\mcses[1]{\ensuremath{\mathit{MCSs}(#1)}\xspace}

%%%%%%%
\newcommand\onestep{\ensuremath{\call{Explain-One-Step}}\xspace}
\newcommand\parallelexplain{\ensuremath{\call{Parallel-One-Step-explain}}\xspace}
\newcommand\onestepo{\ensuremath{\call{Explain-One-Step-ocus}} \xspace}
\newcommand\iterativeonestep{\ensuremath{\call{Explain-One-Step-Iterative}}\xspace}
\newcommand\iterativeonestepDisj{\ensuremath{\call{Explain-One-Step-Iterative-DisjointMCSes}}\xspace}
%%%%%%


\begin{document}
	
%%%%%%% ALGORITHM KEYWORDS
\SetKwInOut{Input}{Input}
\SetKwInOut{OptInput}{Optional}
\SetKwInOut{Output}{Output}
\SetKwInOut{State}{State}
\SetKwInOut{Ext.}{Ext}
\SetKwComment{command}{/*}{*/}
\SetKw{Break}{break}
%%%%%%%

\title{Efficiently Explaining Constraint Satisfaction Problems with Optimal Unsatisfiable Subsets: Extended use-cases and improved algorithms}

\author{}
%\author{\name Emilio Gamba \email emilio.gamba@vub.be \\
%\name Bart Bogaerts \email bart.bogaerts@vub.be \\
%\addr Vrije Universiteit Brussel, Pleinlaan 5\\
%1050 Elsene, Belgium
%\AND
%\name Tias Guns \email tias.guns@vub.be \\
%\addr KULeuven, Oude Markt 13 - bus 5500,\\
%3000 Leuven, Belgium
%}

% For research notes, remove the comment character in the line below.
\researchnote

\maketitle


\section{Example}
The following example will be used as a running example to illustrate these and following concepts.

\begin{example}
Let \emph{C} be a CNF formula over variables $x_1, x_2, x_3$ with the following four clauses:
\[ c_1 := \lnot x_1 \vee \lnot x_2 \vee x_3 \qquad  c_2 := \lnot x_1 \vee  x_2 \vee x_3\] \[  c_3 := x_1 \qquad c_4 := \lnot x_2 \vee \lnot x_3 \]
		
\begin{itemize}
	\item $\Iend$, the final interpretation, is $\{x_1, \lnot x_2,  x_3\}$
	\item $\m{I}$, the current interpretation, is $\{ \emptyset \}$
	\item $\overline{\Iend\setminus I} =  \{\lnot x_1, x_2, \lnot x_3\}$.
\end{itemize}
		
Clause weights \[w_1 = 60, w_2=60, w_3=100, w_4=100, w_5=\infty, w_6=\infty, w_7=\infty, w_8=1, w_9=1, w_{10}=1\]

\begin{itemize}                                      
\item $p_1 := $\texttt{exactly-one}$(\overline{\Iend\setminus \m{I}}) \triangleq $\textit{exactly-one$(\{\lnot x_1, x_2, \lnot x_3 \})$}
\item $p_2 := $\texttt{none-of}$(\Iend\setminus \m{I}) \triangleq $ \texttt{none-of}($\{x_1, \lnot x_2,  x_3\}$)
\end{itemize}


\begin{center}
$\formula :=  $ \emph{C} $ \wedge\  \m{I} \wedge (\overline{\Iend\setminus \m{I}}) \wedge (\Iend\setminus \m{I})$\\
\end{center} 
\begin{center}
\comus$(\formula, f,  \{ p_1 \wedge p_2\})$
\end{center} 
\end{example}


\clearpage

\begin{algorithm}[!h]
\DontPrintSemicolon
		
\caption{$\onestepo(\formulac,f,I,\Iend)$}
\label{alg:oneStepOCUS}
$p \leftarrow$ exactly one of $\overline{\Iend\setminus I}$\;
\Return{$\comus(\formulac\land I\land \overline{\Iend\setminus I}, f, p)$} 
\end{algorithm}
\begin{algorithm}[!h]
\DontPrintSemicolon
$\setstohit  \gets \emptyset$ \; %\label{omus-line1} 
\While{true}{
	$\m{S} \gets \cohs(\setstohit,f,p) $  \;
	\If{ $\lnot \sat(\m{S})$}{\label{alg:ocus-sat-check}
		\Return{$\m{S}$} \;
	}
	$\m{S} \gets  \grow(\m{S},\F) $ \label{alg:ocus:grow}\;
	$\setstohit  \gets \setstohit  \cup \{  \formula \setminus \m{S}\}$  \label{alg:ocus:complement}\;
}
\caption{$\comus(\formula,f,p)$ }
\label{alg:comus}
\end{algorithm}


\begin{algorithm}[!h]
\DontPrintSemicolon
$\setstohit  \gets \emptyset$ \; %\label{omus-line1} 
\While{true}{
	$\m{S} \gets \cohs(\setstohit,f,p) $  \;
	\If{ $\lnot \sat(\m{S})$}{\label{alg:ocus-sat-check}
		\Return{$\m{S}$} \;
	}
	$\m{K} \gets \corrsubsets(\formula, \subsetT)$\;
	$\setstohit  \gets \setstohit  \cup \m{K} $  \label{alg:ocus:complement}\;
}
\caption{$\comus(\formula,f,p)$ }
\label{alg:comus}
\end{algorithm}


\begin{algorithm}[!h]
\DontPrintSemicolon
\Output{A collection of correction subsets}
$\m{C} \gets  \formula \setminus \grow(\m{S}', \formula)$ \;
\Return $ \{ \m{C}\}$
\caption{$\corrsubsets(\formula, \subsetT)$ }
\label{alg:disjmcs}
\end{algorithm}

\begin{algorithm}[!h]
\DontPrintSemicolon
\Output{$\m{K}$: Collection of correction subsets}
$\m{K}, \ \m{S}' \gets \emptyset, \ \m{S}$ \label{alg:disjmcs:init} \;
		
\While{$\sat(\m{S}')$}{
	$\m{C} \gets  \formula \setminus \grow(\m{S}', \formula)$ \;
	$ \m{S}' \gets \m{S}'  \cup  \m{C} $\;
	$ \m{K} \gets \m{K} \cup \m{C} $\;
}
\Return $\m{K}$
\caption{$\corrsubsets(\formula, \subsetT)$ }
\label{alg:disjmcs}
\end{algorithm}

\begin{algorithm}[!h]
\DontPrintSemicolon
\Output{$\m{K}$: Collection of correction subsets}
$\m{K}, \ \m{S}' \gets \emptyset, \ \m{S}$ \label{alg:disjmcs:init} \;
\textit{\texttt{improved}} $\gets true $ \;
\While{$sat(\m{S}') $ and \texttt{improved}}{
	\textit{\texttt{improved}} $\gets false $\;
	$\m{C} \gets  \formula \setminus \grow(\m{S}', \formula)$ \;
	$ \m{K} \gets \m{K} \cup \m{C} $\;

	\For{$c \in \m{C}$}{
		\If{$p( \m{S}'  \cup  \{c\})$}{
			$ \m{S}' \gets \m{S}'  \cup  \{c\} $\;
			\texttt{improved} $\gets true $\,
		}	
	}
}
%\Return $\m{K}, \m{S}'$
\Return $\m{K}$
\caption{$\corrsubsets(\formula, \subsetT, f, p)$ }
\label{alg:disjmcs}
\end{algorithm}

\begin{algorithm}[!h]
\DontPrintSemicolon
\Output{$\m{K}$: Collection of correction subsets}
$\m{K} \gets \emptyset$  \;
$\m{S}'\gets \grow(\m{S}', \formula) $ \;
%
\While{$\m{S}' \neq \m{S}$}{

	$\m{C} \gets  \formula \setminus \m{S}'$ \;
	$ \m{K} \gets \m{K} \cup \m{C} $\;
	$\m{S} \gets \m{S}'$ \tcp{Take a copy}
	\For{$c \in \m{C}$}{
		\uIf{$p( \m{S}'  \cup  \{c\})$}{
			$ \m{S}' \gets \m{S}'  \cup  \{c\} $\;
		}	
	}
		\If{$\sat(\m{S}')$}{
		break\;
		}
	$\m{S}'\gets greedy-\grow(\m{S}', \formula)$ \;

}
%\Return $\m{K}, \m{S}'$
\Return $\m{K}$
\caption{$\corrsubsetsgreedy(\formula, \subsetT, f, p)$ }
\label{alg:disjmcs}
\end{algorithm}

\begin{table}[!h]
\centering
%	\begin{adjustbox}{max width=\columnwidth}
\begin{tabular}{lccccc} 
	%			\bottomrule[2pt]
	\multicolumn{5}{c}{Step 1. Explanation of $x_1$: $c_3 \implies x_1$}\\
	\toprule[2pt]
	  & $\m{S}$                             & $\sat(\m{S})$    & \grow($\m{S}$, $\formula$)                                 & $\m{C} \gets  \formula \setminus \grow(\m{S}, \formula)$ &  $\setstohit  \gets \setstohit  \cup \{  \m{C}\}$ \\ 
	\toprule[2pt]
	1 & $\{ \lnot x_1\}$                    & $\mathit{true}$  & $\{c_1, c_2, c_4, \lnot x_1, \lnot x_2, x_3 \}$            & $\{ c_3, x_1, x_2, \lnot x_3 \}$ &  $\{ \{c_3, x_1, x_2, \lnot x_3 \}\}$ \\
	\midrule	
	2 & $\{ \lnot x_1, c_3, x_1, x_2, \lnot x_3\}$  & $\mathit{false}$  &  & &\\
	%			\bottomrule
				
\end{tabular}
%	\end{adjustbox}
\caption{Sequence of explanation generated with \onestepo and \corrsubsets}
\label{tab:explanation-steps-expanded:step1}
\end{table}




\begin{table}[!h]
\centering
%	\begin{adjustbox}{max width=\columnwidth}
\begin{tabular}{lcccc} 
	%			\bottomrule[2pt]
	\multicolumn{5}{c}{Step 1. Explanation of $x_1$: $c_3 \implies x_1$}\\
	\toprule[2pt]
	  & $\m{S}$                             & $\sat(\m{S})$    & \grow($\m{S}$, $\formula$)                                 & $\setstohit  \gets \setstohit  \cup \{  \formula \setminus \m{S}\}$ \\ 
	\toprule[2pt]
	1 & $\{ \lnot x_1\}$                    & $\mathit{true}$  & $\{c_1, c_2, c_4, \lnot x_1, \lnot x_2, x_3 \}$            & $\{ \{c_3, x_1, x_2, \lnot x_3 \}\}$                                \\
	\midrule	
	2 & $\{\lnot x_3\}$                     & $\mathit{true}$  & $\{c_1 , c_2 , c_4, \lnot x_1 , \lnot x_2 ,  \lnot x_3 \}$ & $\{ ..., \{c_3, x_1 , x_2 , x_3 \}\}$                               \\  
	\midrule
	3 & $\{x_2\}$                           & $\mathit{true}$  & $\{c_1 , c_2 , c_3, x_1 , x_2 , x_3 \}$                    & $\{ ..., \{c_4, \lnot x_1 , \lnot x_2 , \lnot x_3\} \}$             \\
	\midrule
	4 & $\{x_2, c_4\}$                      & $\mathit{true}$  & $\{  c_1 , c_2 , c_4, \lnot x_1 , x_2 , \lnot x_3\}$       & $\{ ..., \{x_1 , \lnot x_2 , x_3 , c_3 \}\}$                        \\
	\midrule
	5 & $\{\lnot x_3, c_3\}$                & $\mathit{true}$  & $\{   c_1  , c_3 , c_4, x_1, \lnot x_2 , \lnot x_3 \}$     & $\{ ...,  \{ c_2, \lnot x_1 , x_2 , x_3\}\}$                        \\
	\midrule
	6 & $\{\lnot x_1, c_3\}$                & $\mathit{false}$ & -                                                          & -                                                                   \\
	%			\bottomrule
	\\
	\multicolumn{5}{c}{Step 2. Explanation of $x_3$: $c_1 \wedge c_2 \wedge x_1 \implies x_3$}\\
	\toprule[2pt]
	7 & $\{x_1 , \lnot x_3 , c_2\}$         & $\mathit{true}$  & $\{   c_2 , c_3 , c_4,x_1 , x_2 , \lnot x_3 \}$            & $\{ ..., \{c_1, \lnot x_2 , x_3 \} \}$                              \\
	\midrule
	%		
	%		7 & $\{ x_1 , \lnot x_3 , c_2\}$ & $\mathit{true}$ & $\{ c_2 , c_3 , c_4,x_1 , x_2 , \lnot x_3 \}$ & $\{ ..., \{c1, \lnot x2 , x3 \}\}$  \\ \midrule
	8 & $\{ x_1 , \lnot x_3 , c_1 , c_2 \}$ & $\mathit{false}$ & -                                                          & -                                                                   \\
	%			\bottomrule
	\\
	\multicolumn{5}{c}{Step 3. Explanation of $\lnot x_2$: $c_4 \wedge x_3 \implies \lnot x_2$}\\
	\toprule[2pt]
	9 & $\{ x_2 , x_3, c_4 \}$              & $\mathit{false}$ & -                                                          & -                                                                   \\
					
\end{tabular}
%	\end{adjustbox}
\caption{Sequence of explanation generated with \onestepo}
\label{tab:explanation-steps-expanded:step1}
\end{table}




\clearpage
\section{Iterative searching for \omus}

\begin{algorithm}[!h]
\DontPrintSemicolon
\tcp{Initilisation of Heap with tuples $\langle $cost$, $ literal$, $ hs$ \rangle$}
$lb \gets$ \texttt{heap()}\;
\ForEach{$\ell \in  (\Iend \setminus I)$}{
	$	\formula_{\ell} \gets  \formulac \wedge (\Iend \setminus \ell) \wedge \{\lnot \ell \}$ \;
	$ \m{S} \gets \grow(\emptyset, \ \formula_{\ell})$\; 
	$\setstohit_{\ell}  \gets \setstohit_{\ell}  \cup \{  \formula_{\ell} \setminus \m{S}\}$ \;
	\texttt{lb.push($\langle 0, \ell, \emptyset \rangle $)}\;
}
\tcp{\comus loop}
\While{$true$}{
				
	$(cost, \ell, \m{S}) \gets$\texttt{lb.pop()} \tcp*{Literal with lowest bound}
				
	\If{ $\lnot \sat(\m{S})$}{
		\Return{$\m{S}$} \tcp*{\color{Green}OCUS found}
	}
				
	$\m{S} \gets  \grow(\m{S}, 	\formula_{\ell})$ \tcp*{Hitting set grown to a Satisfiable Subset}
				
	$\setstohit_{\ell}  \gets \setstohit_{\ell}  \cup \{  \formula_{\ell} \setminus \m{S}\}$ \tcp*{Correction Subset added to sets-to-hit}
				
	$\m{S} \gets \cohs(\setstohit_{\ell},f,p)$ \tcp*{Optimal hitting set}
				
	\texttt{lb.push($\langle f(\m{S}), \ell, \m{S} \rangle$)}\;
}
\caption{$\iterativeonestep(\formulac,f,I, \Iend)$}
\label{alg:iterative_ocus}
\end{algorithm}
\clearpage


\clearpage
\begin{algorithm}[!h]
\DontPrintSemicolon
\Input{$\textbf{p}$: constraint(s) on \comus}
\Output{$\m{K}$: Disjoint correction subsets}
\Output{$\m{S}'$: Unsatisfiable Subset of $\m{F}$ satisfying $\textbf{p}$}
$\m{K}, \ \m{S}' \gets \emptyset, \ \m{S}$ \label{alg:disjmcs:init} \;
		
\While{$\sat(\m{S}')$}{
	$\m{C} \gets  \formula \setminus \grow(\m{S}', \formula)$ \;
	$ \m{K} \gets \m{K} \cup \m{C} $\;
	\tcp{`Greedy' Extending the optimal constrained hitting set when adding the elements of the complement to a hitting set}
	\For{$c \in \m{C}$}{
		\uIf{$p_1(c \wedge \m{S'}) \wedge .. \wedge p_n(c \wedge \m{S'})$ }{
			$S' \gets S' \cup {c}$\;
		}
	}
	\uIf{$f(\m{S}') > ub $}{
		\Return $\m{K}, \m{S}'$
	}
	

}
\Return $\m{K}, \m{S}'$
\caption{$\disjcorrconstrained(\formula, f, \textbf{p} ,  \subsetT, ub=\infty)$ }
\label{alg:disjmcs}
\end{algorithm}
\clearpage

\subsection{Explanations with Bounded OUS.}
%Context: how to do OUS 

\clearpage

\begin{algorithm}[!h]
\DontPrintSemicolon
\tcp{Initilisation of Heap with tuples $\langle $lb$,$ ub $, $ literal$, $ hs$ \rangle$}
$lb \gets$ \texttt{heap()}\;

\ForEach{$\ell \in  (\Iend \setminus I)$}{
	$	\formula_{\ell} \gets  \formulac \wedge (\Iend \setminus \ell) \wedge \{\lnot \ell \}$ \;
	$p := \{\ $\texttt{use-only}$(\{ I \wedge (\lnot \ell) \}), \{ \lnot \ell\} \ \}$\;
	$ \m{K}, \m{S} \gets  \disjcorrconstrained(\formula_{\ell} , \emptyset, p, ub=\infty) $\;
	$ \setstohit_{\ell} \gets \setstohit_{\ell} \cup \m{K}$\;
	\uIf{$p(\m{S})$}{
		$ub \gets f(\m{S})$\;
		$\m{S}_{best} \gets \m{S}$\;
	}
	$\m{S} \gets \cohs(\setstohit_{\ell},f,p)$\;
	\texttt{lb.push($\langle \ f(\m{S}), \ ub,\ \ell,\ \m{S} \rangle$)}\;
	%		\texttt{lb.push($\langle 0, \ell, \emptyset \rangle $)}\;
}
\tcp{\comus loop}
\While{$lb \neq ub$}{
				
	$( \ lb, \  ub,  \ \ell,  \ \m{S}) \gets$\texttt{lb.pop()} 
				
	\If{ $\lnot \sat(\m{S})$}{
		\Return{$\m{S}$} \tcp*{\color{Green}OCUS found}
	}
				
	%		$\m{S} \gets  \grow(\m{S}, 	\formula_{\ell})$ \tcp*{Hitting set grown to a Satisfiable Subset}
	$p := \{$\texttt{use-only}$(\{ I \wedge (\lnot \ell) \})\}$\;
	$ \m{K}, \m{S}' \gets  \disjcorrconstrained(\formula_{\ell}, \m{S}, p, ub)  $\;
	\uIf{$f(\m{S}') < ub$}{
		$ub \gets f(\m{S}')$\;
		$\m{S}_{best} \gets \m{S}$\;
	}
	$ \setstohit_{\ell} \gets \setstohit_{\ell} \cup \m{K}$\;
%	$\setstohit_{\ell}  \gets \setstohit_{\ell}  \cup \disjcorrconstrained(\formula_{\ell}, \m{S}, p, ub) $\;
	$\m{S} \gets \cohs(\setstohit_{\ell},f,p)$ \tcp*{Optimal hitting set}
	\texttt{lb.push($\langle f(\m{S}), ub, \ell, \m{S} \rangle$)}\;
}
\Return{$\m{S}_{best}$}
\caption{$\iterativeonestepDisj(\formulac,f,I, \Iend)$}
\label{alg:iterative_ocus}
\end{algorithm}
\clearpage



Algorithm \ref{alg:comus} computes an \comus for a given $\formula$, cost function $f$, and predicate $p$. 


\section{Experimental Results}\label{sec:experiments}
\subsection{Datasets}
\begin{itemize}
\item Logic Grid puzzles not enough
\item Sudoku puzzles of varying difficulty (`simple', `easy', `intermediate', `expert') generated using qqwing.
\item Instances from Demystify 
\end{itemize}

\subsection{RQ1. Cost comparison between MUS and Constrained for Logic Grid Puzzles+ Extra Demystify puzzles}



\subsection{RQ2. Domain-specific vs disj. mcs grow}

\subsection{RQ3. Effect of COnstrainedness on the time required to computed explanations}

\subsection{RQ4. Does re-use of information across the different iterations improve efficiency?}

\subsection{RQ5. Does appropriate sorting of literals lead to faster generation of individual explanations and explanation sequences ?}



% Graph of most important configurations
% 1) improvement with disjoint mcs enumeration
% 2) Replacing the grow with other strategies: DIJS MCSes
% 3) Replacing greedy loop with iterative search
% 4) Constrained

%\subsection{Hitting set computation}
% someway to present set coverage between hittingsets
% \number of hitting sets computed for the most important configuration
% 


\section{Conclusion and Future work}

\vskip 0.2in
\bibliography{krrlib,ref}
\bibliographystyle{theapa}

\end{document}





